\chapter{Kết luận và hướng phát triển}
\label{chap:conclusion}

Chương này tổng kết lại những kết quả đạt được của dự án nghiên cứu và phát triển hệ thống Self-Evolving AI Agents trên nền tảng NIKA, đồng thời thảo luận thẳng thắn về những hạn chế hiện tại và đề xuất định hướng phát triển trong tương lai.

\section{Kết quả đạt được của dự án}
\label{sec:accomplishments}

Sau quá trình nghiên cứu và thực nghiệm nghiêm túc, dự án đã đạt được một số kết quả quan trọng sau:
\begin{itemize}
    \item **Nghiên cứu lý thuyết sâu rộng:** Tổng hợp và hệ thống hóa các nghiên cứu hàng đầu về AI Agent tự tiến hóa bao gồm cả hướng tiến hóa bộ nhớ (MemInsight, LightMem, A-Mem) và tiến hóa công cụ (DRAFT, TTE, Alita).
    \item **Xây dựng kiến trúc hệ thống hoàn chỉnh:** Tích hợp thành công AI Agent với môi trường giả lập mạng Kathará thông qua giao thức Model Context Protocol (MCP) chuẩn hóa. Hệ thống vận hành tốt trên 3 luồng chẩn đoán phổ biến (ReAct, Plan \& Execute, Reflexion).
    \item **Hiện thực hóa thành công các mô-đun tiến hóa:**
        \begin{enumerate}
            \item Mô-đun Bộ nhớ quy trình tự tiến hóa (Composable Procedural Memory Module) sử dụng cấu trúc lai PostgreSQL/Qdrant hỗ trợ nén vết chẩn đoán và cập nhật ngoại tuyến, đặc biệt là cơ chế score-gated validation để tránh rò rỉ đáp án của benchmark.
            \item Mô-đun Tiến hóa công cụ (Tool Evolution Module) cho phép tự Mastery tài liệu hướng dẫn và sinh công cụ mới (composite workflows và generated Python helpers) có cơ chế kiểm tra AST và thực thi Docker sandbox an toàn.
        \end{enumerate}
    \item **Đánh giá thực nghiệm chi tiết:** Kiểm chứng hệ thống trên 30 incident của benchmark NIKA. Kết quả cho thấy Evolved Agent đạt điểm F1 trung bình vượt trội so với baseline (tăng từ 0.0111 lên 0.1611) và giải quyết thành công thêm nhiều sự cố phức tạp liên quan đến mất cấu hình IP và lỗi định tuyến BGP.
\end{itemize}

\section{Những hạn chế hiện tại}
\label{sec:limitations}

Mặc dù đạt được những kết quả khả quan, hệ thống vẫn tồn tại một số hạn chế cần khắc phục:
\begin{itemize}
    \item **Hiện tượng quá khớp (Overfitting) trên phân đoạn Transfer:** Các công cụ và bài học được tiến hóa trên kịch bản mạng cũ đôi khi bị quá khớp với cấu trúc hoặc giao thức đó. Khi chuyển giao sang phân đoạn \texttt{transfer} với các lỗi định tuyến động mới (OSPF), agent dễ gặp lỗi gọi công cụ do không tương thích tham số, làm giảm F1 score.
    \item **Tăng chi phí Token:** Việc tải thêm các mô tả Mastery dài và các gợi ý bài học từ bộ nhớ vào ngữ cảnh prompt làm tăng lượng token trung bình tiêu thụ mỗi ca từ khoảng 64,468 lên 80,771 tokens.
    \item **Rủi ro an toàn thực thi ở test-time:** Mặc dù đã có cơ chế kiểm tra cú pháp AST và thực thi Docker sandbox cô lập hoàn toàn mạng để chạy thử nghiệm các công cụ Python mới sinh, việc cho phép LLM tự động sinh mã nguồn thực thi thời gian chạy vẫn tiềm ẩn rủi ro bảo mật lớn nếu đưa vào hệ thống mạng viễn thông thực tế.
\end{itemize}

\section{Định hướng phát triển tương lai}
\label{sec:future_work}

Từ những kết quả và hạn chế trên, dự án định hướng một số nội dung nghiên cứu tiếp theo bao gồm:
\begin{itemize}
    \item **Cải thiện khả năng tổng quát hóa (Generalization):** Nghiên cứu các cơ chế trừu tượng hóa tham số mạnh mẽ hơn trong lúc chẩn đoán và distill công cụ, giúp các công cụ composite workflow dễ dàng thích ứng với các cấu trúc mạng và giao thức khác nhau (giải quyết bài toán Transfer).
    \item **Tích hợp cơ chế Human-in-the-loop:** Bổ sung giao diện cho phép các kỹ sư mạng con người tham gia phản hồi, đánh giá và phê duyệt (approve) các công cụ/quy trình mới sinh ra cũng như các ghi nhớ quan trọng trước khi ghi nhận chính thức vào state.json.
    \item **Tối ưu hóa hiệu năng và chi phí:** Tinh chỉnh các mô hình ngôn ngữ lớn chuyên biệt cho miền mạng (Domain-specific LLMs) để giảm kích thước prompt chẩn đoán và tối ưu hóa thời gian thực thi các Kathará labs.
    \item **Mở rộng kịch bản sự cố:** Đánh giá agent tự tiến hóa trên các kịch bản mạng có quy mô lớn hơn (m, l) và các sự cố phức tạp hơn (compose nhiều lỗi trong một incident).
\end{itemize}
